%% kapa: exam=1 round=36 year=2025 period=2 form=A subject=accounting subject_ko=회계학 seq=19 number=59
%% source: Q-Net 감정평가사 기출문제 (cst00309) artlSeq=5231525 file=raw/1008/5231525/2025년 제36회 감정평가사 1차시험 2교시 문제지 원본.pdf
%% conversion: kordoc 4.24.51 -> tools/gen1.py (auto); review=proofread-sonnet; flags=none
\begin{question}{59}
충당부채, 우발부채 및 우발자산에 관한 설명으로 옳은 것은?
\begin{choices}
  \item 미래의 예상 영업손실은 충당부채로 인식하여야 한다.
  \item 손실부담계약을 체결하고 있는 경우에는 관련된 현재의무를 충당부채로 인식하지 않는다.
  \item 보수주의 관점에서 우발자산은 재무제표에 인식하지 아니하나, 우발부채는 재무제표에 인식한다.
  \item 충당부채와 관련하여 포괄손익계산서에 인식한 비용은 제삼자의 변제와 관련하여 인식한 금액과 상계하여 표시할 수 있다.
  \item 예상되는 자산 처분이 충당부채를 생기게 한 사건과 밀접하게 관련되었다면, 예상되는 자산 처분이익은 충당부채를 측정하는 데 고려하여야 한다.
\end{choices}
\end{question}
\answer{4}
